\section{Hermitian forms and quadratic forms}

In all that follows in this Chapter, unless expressly stated otherwise, A denotes a ring and E a left A-module. One supposes A endowed with an involutive antiautomorphism J, denoted \(\alpha \to \overline{\alpha}\) ; one therefore has \((\overline{\alpha + \beta}) = \overline{\alpha} + \overline{\beta}\) , \((\overline{\alpha\beta}) = \overline{\beta}.\overline{\alpha}\) and \(\overline{\overline{\alpha}} = \alpha\) for all \(\alpha\) , \(\beta\) in A. Unless expressly stated otherwise, the sesquilinear forms considered are right sesquilinear ( \(\S 1\) , no. 2, def. 4) for this antiautomorphism.

